\worksheettitle
  {Одинаковые цифры --- разные значения}
  {\modulemark{M01\(\star\)} Усложнение}

\studentnameblock

\begin{problembox}[Исследование]
  Найдите число \(7\,707\). В его записи три цифры \(7\).

  \begin{enumerate}[label=\textbf{\arabic*.}]
    \item Назовите значение каждой цифры \(7\):
      \answerblank[25mm], \answerblank[25mm], \answerblank[25mm].
    \item Во сколько раз значение первой цифры \(7\) больше значения
      последней? \answerblank[20mm].
    \item Переставьте одну цифру \(7\), чтобы получилось другое число, и
      объясните, какое значение изменилось.
      \writinglines{2}
  \end{enumerate}
\end{problembox}

\begin{practicebox}[Создайте свой пример]
  Запишите пятизначное число, в котором одна и та же ненулевая цифра
  встречается не менее трёх раз.

  Число: \answerblank[35mm].

  Назовите разряды и значения повторяющейся цифры:
  \writinglines{3}
\end{practicebox}

\begin{reflectionbox}[Обобщение]
  Закончите предложение: «Чтобы узнать значение цифры, недостаточно увидеть
  саму цифру, потому что…»
  \writinglines{1}
\end{reflectionbox}
